\documentclass[11pt]{article}
\usepackage[margin=1in]{geometry}
\usepackage{amsmath,amssymb,amsthm,hyperref}
\emergencystretch=2em
\newtheorem{theorem}{Theorem}
\title{Conjecture 00000002819: independent additive Gaussian sums have no affine large-deviation rate}
\author{GitHub: gaochengzhecpu}
\date{}
\begin{document}
\maketitle
\begin{abstract}
Independent standard Gaussian increments give an additive, hence
subadditive, process. Its sample-mean laws admit no affine real-valued
large-deviation rate function at the usual linear speed. We prove this
directly from genuine probability bounds and the full large-deviation
inequalities. The obstruction occurs inside the support, which is all
of the real line. The actual Gaussian rate is quadratic.
\end{abstract}

\section*{Model and reading of the conjecture}
The source asserts that, in the independent case, the large-deviation
rate of subadditive functionals is always a linear function. We use
the standard meaning of a rate function $I(x)$ as a function of the
deviation level $x$. We even allow an affine rate $ax+b$; excluding
these excludes homogeneous linear rates as well. This is distinct from
the standard \emph{speed} $N$ in $\log\mathbb P(\text{deviation})
\sim-NI(x)$. The latter can be linear in $N$ when $I(x)$ is nonlinear.
If the word ``linear'' was intended to describe only the speed, this
is a different statement from linearity of the rate function and is
not refuted here.

On the actual probability space
\[
 \Omega=\mathbb R^{\mathbb N},\qquad
 \mathbb P=\bigotimes_{j\geq0} N(0,1),
\]
let $X_j(\omega)=\omega_j$. These coordinates are independent and
identically distributed standard Gaussians. Define
\[
 A_{m,n}=\sum_{j=0}^{n-1}X_{m+j},\qquad
 S_N=A_{0,N},\qquad M_N=S_N/N\quad(N\geq1).
\]
For every sample, $A_{m,n+k}=A_{m,n}+A_{m+n,k}$, so these functionals
are genuinely additive and satisfy the subadditive inequality with
equality. In particular, no dependence correction is present.
Write $\nu_N$ for the probability law of $M_N$.

\section*{The standard LDP and its exponential form}
For a finite real-valued nonnegative lower semicontinuous rate $I$,
the full LDP at speed $N$ consists of
\[
 \begin{aligned}
 \liminf_{N\to\infty}N^{-1}\log\nu_N(G)&\geq-\inf_{x\in G}I(x)
                  &&\text{for every open }G,\\
 \limsup_{N\to\infty}N^{-1}\log\nu_N(F)&\leq-\inf_{x\in F}I(x)
                  &&\text{for every closed }F,
 \end{aligned}
\]
where $\log0=-\infty$ and $\inf\varnothing=+\infty$.
The Lean definition uses the equivalent epsilon exponential formulation:
for every open $G$, every $x\in G$, and every $\varepsilon>0$,
eventually
\[
 e^{-N(I(x)+\varepsilon)}\leq\nu_N(G),
\]
and for every closed $F$, every real uniform lower bound $a$ for
$I$ on $F$, and every $\varepsilon>0$, eventually
\[
 \nu_N(F)\leq e^{-N(a-\varepsilon)}.
\]
These conditions range over \emph{all} open and closed sets, not just
the sets used in the contradiction. Their equivalence to the displayed
LDP follows by taking logarithms and then taking the supremum over
$x$ in the lower bound and over uniform lower bounds $a$ in the upper
bound. Conversely, the logarithmic inequalities give each eventual
epsilon bound. Empty sets satisfy the corresponding bounds automatically.
Finite real-valued rates suffice for the proposed affine rates.

\begin{theorem}
For the independent additive Gaussian process above, there are no
$a,b\in\mathbb R$ such that $I(x)=ax+b$ is an LDP rate for $(\nu_N)$
at speed $N$. Every nonempty real open interval has positive
$\nu_N$-probability for every $N\geq1$.
\end{theorem}
\begin{proof}
The genuine moment generating function of $S_N$ is
\[
 \mathbb E e^{tS_N}
 =\prod_{j=0}^{N-1}\mathbb E e^{tX_j}
 =e^{Nt^2/2}.
\]
Chernoff's inequality with $t=1$ gives
\[
 \nu_N((1,\infty))\leq\mathbb P(S_N\geq N)
 \leq e^{-N}\mathbb E e^{S_N}=e^{-N/2}.                 \tag{1}
\]
Suppose $I(x)=ax+b$ were a rate. Nonnegativity on the entire real
line forces $a=0$: if $a\ne0$, substituting $x=(-b-1)/a$ gives
$I(x)=-1$. Thus $I=b\geq0$. If $b>0$, the closed-set upper bound
for $F=\mathbb R$, with $\varepsilon=b/2$, would give
$1\leq e^{-Nb/2}<1$ for all sufficiently large $N$. Hence $b=0$.
The open-set lower bound for $G=(1,2)$, at $x=3/2$ with
$\varepsilon=1/4$, now gives, eventually,
\[
 e^{-N/4}\leq\nu_N((1,2))\leq\nu_N((1,\infty))\leq e^{-N/2}.
\]
This is impossible for every positive $N$.

Finally, if $u<v$, the event that all of $X_0,\ldots,X_{N-1}$ lie
in $(u,v)$ has probability
$\mathbb P(X_0\in(u,v))^N>0$ and forces $M_N\in(u,v)$.
The strict positivity follows from the positive Gaussian density.
Thus the argument uses an actual open interval inside the full real
support, not an interval outside the possible values.
\end{proof}

\section*{The quadratic rate, for context}
The actual increment cumulant generating function is $\Lambda(t)=t^2/2$.
Its exact Legendre transform is
\[
 \sup_{t\in\mathbb R}\{tx-\Lambda(t)\}
 =\frac{x^2}{2},
 \qquad tx-\frac{t^2}{2}
 =\frac{x^2}{2}-\frac{(t-x)^2}{2}.
\]
This transform computation, including its non-affinity, is separately
formalized using the actual Gaussian cumulant generating function.

For completeness, the Gaussian LDP itself can also be seen directly
on paper. Gaussian convolution gives the density of $M_N$ as
$\sqrt{N/(2\pi)}e^{-Nx^2/2}$. For any open set and a point $x$ in
it, integrate over a sufficiently small interval about $x$; continuity
of $x^2/2$ gives the LDP lower bound, because the logarithm of the
prefactor divided by $N$ tends to zero. For a nonempty closed set
$F$, put $r=\inf_{x\in F}|x|$. If $r=0$, the probability bound by one
gives the upper bound. If $r>0$, the two Chernoff tail bounds give
$\nu_N(F)\leq2e^{-Nr^2/2}$. This proves the upper bound as well.
The empty closed set is immediate. Thus the example has a real,
nonlinear LDP rate; the contradiction above does not merely exploit
nonexistence of a large-deviation principle.

\section*{Formal verification and scope}
Lean constructs Mathlib's actual infinite product Gaussian probability
measure, proves its coordinate marginals and independence, and proves
the block-sum additive and subadditive identities. It defines the
sample means and their actual pushforward measures. From Gaussian
MGFs and independence it proves the MGF of each partial sum, the
required exponential integrability, and the actual Chernoff bound (1).
It also proves positive probability for every nonempty real open interval.

The full epsilon-form LDP is defined with nonnegativity, lower
semicontinuity, and bounds for every open and every closed set.
Lean negates the existential proposition that an affine rate satisfies
this definition. The concluding theorem combines the genuine independent
Gaussian coordinates, samplewise subadditivity, and that negation:
\begin{center}\small
\texttt{Conjecture2819.independent\_subadditive\_counterexample}.
\end{center}
No surrogate numerical function replaces the random process or the
probability inequalities. The supplementary Cramer-transform computation
is also formalized. The optional paper derivation that the quadratic
rate satisfies the full LDP is not formalized; it is unnecessary for
the formal proof that no affine LDP rate can exist. The source's separate
correlation-length correction is not analyzed. The axiom audit of the
principal theorems lists only \texttt{propext}, \texttt{Classical.choice},
and \texttt{Quot.sound}.

\section*{Source and reproduction}
The unchanged bilingual statement is included as \texttt{SOURCE.md},
from upstream commit
\begin{center}\small
\texttt{b355e42906e6eaf52bc7e9d005e0ef02d8aa7400}.
\end{center}
The public source is
\href{https://github.com/The-Last-Math-Competition/The-Last-Math-Competition/blob/b355e42906e6eaf52bc7e9d005e0ef02d8aa7400/conjectures/00000002819.md}{conjecture 00000002819}.
Use Lean 4.19.0 and Mathlib commit
\begin{center}\small
\texttt{c44e0c8ee63ca166450922a373c7409c5d26b00b}.
\end{center}
In the \texttt{lean} directory run \texttt{lake build}, then
\begin{center}
\texttt{lake env lean -DwarningAsError=true Main.lean}.
\end{center}
This draft was developed by a delegated AI agent with a self-review;
the coordinating agent reviews it separately.
\end{document}
